\documentclass[11pt,reqno]{amsart}

\usepackage[T1]{fontenc}
\usepackage{lmodern}
\usepackage{microtype}
\usepackage{amsmath,amssymb,mathtools}
\usepackage{booktabs}
\usepackage{enumitem}
\usepackage{xcolor}
\usepackage[margin=1.05in]{geometry}
\usepackage[colorlinks=true,
            linkcolor=blue!45!black,
            citecolor=blue!45!black,
            urlcolor=blue!45!black]{hyperref}

\allowdisplaybreaks
\numberwithin{equation}{section}

\newtheorem{theorem}{Theorem}[section]
\newtheorem{proposition}[theorem]{Proposition}
\newtheorem{lemma}[theorem]{Lemma}
\newtheorem{corollary}[theorem]{Corollary}
\newtheorem{conjecture}[theorem]{Conjecture}
\newtheorem{problem}[theorem]{Problem}
\theoremstyle{definition}
\newtheorem{definition}[theorem]{Definition}
\newtheorem{construction}[theorem]{Construction}
\theoremstyle{remark}
\newtheorem{remark}[theorem]{Remark}

\newcommand{\N}{\mathbb N}
\newcommand{\Z}{\mathbb Z}
\newcommand{\Q}{\mathbb Q}
\newcommand{\R}{\mathbb R}
\newcommand{\Ztwo}{\mathbb Z_2}
\newcommand{\eps}{\varepsilon}
\newcommand{\vtwo}{v_2}
\newcommand{\vthree}{v_3}
\newcommand{\one}{\mathbf 1}
\newcommand{\abs}[1]{\left|#1\right|}
\newcommand{\set}[1]{\left\{#1\right\}}
\newcommand{\floor}[1]{\left\lfloor #1\right\rfloor}
\newcommand{\ceil}[1]{\left\lceil #1\right\rceil}
\newcommand{\angles}[1]{\left\langle #1\right\rangle}
\DeclareMathOperator{\wt}{wt}
\DeclareMathOperator{\Col}{Col}

\title[Stationary inverse-fiber extinction]
{Stationary Inverse-Fiber Extinction and the Synchronization Obstruction\\
in the Collatz Problem}

\author{Anton Shakov}
\address{Principia Math}
\date{August 2026}

\subjclass[2020]{11B83, 37P99, 11J87, 68Q42}
\keywords{Collatz conjecture, paradoxical sequence, inverse fiber, parity vector,
$2$-adic conjugacy, $3$-adic carry, Christoffel word, min-plus dynamics}

\begin{document}

\begin{abstract}
We combine the paradoxical-sequence framework of Rozier and Terracol with
exact affine, $2$-adic, $3$-adic, and inverse-tree descriptions of the
Collatz map.  The central exact object is the least positive element
\[
M_r(n)=\min\set{y>0\text{ odd}:U^r(y)=U^r(n)}
\]
of the $r$-step inverse fiber above the actual accelerated endpoint of an
odd integer $n$.  We prove that the Collatz conjecture is equivalent to the
finite certificate condition $M_r(n)<n$ for some $r$ and every odd $n>1$.
We then derive a finite-support characterization of positive integers in
parity coordinates, a mixed dyadic--ternary corner criterion, a quantitative
resonance theorem for primitive paradoxical blocks, an exact winding formula
for adjacent parity transpositions, a barrier-preserving reduction to a
mechanical extremal word, and an equal-slope carry criterion producing a
smaller coalescing ancestor.

These results isolate the remaining obstruction as a synchronization problem:
a single positive ordinary integer would have to remain the least positive
point in every finite inverse fiber while its left $2$-adic seed data, right
$3$-adic endpoint data, and Archimedean order data stay compatible at all
scales.  We formulate ten concrete research programs aimed at this
obstruction, drawing on continued fractions, Ostrowski numeration, min-plus
geometry, moving-target Diophantine approximation, $p$-adic interpolation,
fractal uncertainty, autoconjugacy, and substitution cocycles.  No
unconditional proof of the Collatz conjecture is claimed.  The final theorem
proves Collatz conditional on one sharply stated stationary inverse-fiber
extinction assertion, and also gives a two-lemma completion architecture
through first-crossing descent and high-surplus carry innovation.
\end{abstract}

\maketitle

\section{Introduction}

Let
\begin{equation}\label{eq:shortcut-map}
T(n)=
\begin{cases}
 n/2,&n\equiv0\pmod2,\\[1mm]
 (3n+1)/2,&n\equiv1\pmod2
\end{cases}
\end{equation}
be the shortcut Collatz map.  The Collatz conjecture asserts that every
positive integer eventually enters the cycle $1\leftrightarrow2$.  The
conjecture has been computationally verified below $2^{71}$ by Barina
\cite{Barina2025}, while Tao proved that almost every orbit, in logarithmic
density, attains values below any prescribed function tending to infinity
\cite{Tao2022}.  Neither result excludes one exceptional divergent orbit or
one nontrivial positive cycle.

Rozier and Terracol introduced the systematic study of \emph{paradoxical}
finite Collatz segments: blocks whose multiplicative coefficient is below
one even though the endpoint is not below the starting value
\cite{RozierTerracol2025}.  Their work connects the finiteness of such blocks
to the full Collatz problem and to the coefficient-stopping-time conjecture
of Terras \cite{Terras1976}.  Subsequent work has investigated fixed-length
paradoxical vectors \cite{Niu2026}, two-ended $2$-adic/$3$-adic searches
\cite{Kramer2026}, and extremal parity words governed by Christoffel
structures \cite{FernandezIbanez2026}.

The purpose of the present paper is not to claim a proof, but to make the
remaining global obstruction mathematically exact.  At finite depth, the
symbolic dynamics are completely explicit:

\begin{itemize}[leftmargin=2em]
\item every parity or valuation word has an affine formula;
\item every finite parity word determines one residue class modulo a power of
      two;
\item terminal endpoint conditions modulo powers of three are determined by
      terminal valuation suffixes;
\item every competing inverse path can be compared by an exact affine score.
\end{itemize}

The difficulty is the passage from ``every depth is realizable'' to ``one
fixed positive ordinary integer realizes every depth.''  Every finite parity
word occurs on a terminating orbit, so no bounded-prefix obstruction can
separate a hypothetical counterexample from longer and longer terminating
shadows.  The exceptional object would have to be a \emph{stationary inverse
fiber minimum}: the same positive integer would remain the least positive
ancestor of its actual endpoint at every depth.

Our main exact theorem states that excluding such stationary minima is
\emph{equivalent} to Collatz.  This gives the narrowest single completion
target found in the present framework.

\subsection*{Main contributions}

The paper contains the following exact results.

\begin{enumerate}[label=(\roman*),leftmargin=2.5em]
\item We prove that Collatz is equivalent to stationary inverse-fiber
      extinction: for every odd $n>1$, some finite endpoint of the orbit has a
      smaller positive odd ancestor.
\item We give a quantitative resonance theorem for primitive paradoxical
      blocks.  A paradox with a large minimum forces an exceptionally accurate
      upper rational approximation to $\log_2 3$ and nearly maximal affine
      carry.
\item We characterize the positive-integer image of the parity conjugacy by
      finite support of an intrinsic inverse-carry sequence.  Inside the
      coefficient barrier this becomes a mixed $2$-adic/$3$-adic corner
      condition.
\item We prove an adjacent-transposition winding identity.  Moving a parity
      pattern $10$ to $01$ has a monotone real carry effect but a seed change
      containing one binary wrap bit; this wrap bit is the precise obstruction
      to transferring Christoffel carry extremality to descent.
\item We prove that every primitive first-crossing word can be moved, through
      barrier-preserving adjacent transpositions, to a canonical mechanical
      word.
\item We introduce an equal-slope carry spectrum and prove that one small
      positive carry return gives a smaller positive coalescing ancestor.
\item We formulate ten explicit attacks on synchronization and state the exact
      theorem each would have to prove.
\end{enumerate}

\subsection*{Status}

Several statements below are theorems and are proved in full.  The ten
research programs are explicitly marked as open targets.  The Collatz
conjecture remains conditional on stationary inverse-fiber extinction, or on
the more structured pair of assertions called first-crossing descent and
high-surplus carry innovation in Section~\ref{sec:completion}.

\section{Affine word calculus}\label{sec:affine}

\subsection{Shortcut parity words}

Let
\[
w=(\eps_0,\ldots,\eps_{k-1})\in\set{0,1}^k,
\qquad
s_j=\sum_{i=0}^{j-1}\eps_i,
\qquad s=s_k.
\]
Define
\begin{equation}\label{eq:shortcut-carry-rec}
C_0=0,
\qquad
C_{j+1}=3^{\eps_j}C_j+\eps_j2^j.
\end{equation}

\begin{lemma}[Affine parity-word identity]\label{lem:shortcut-affine}
If $w$ is the first $k$ parity bits of $n$, then
\begin{equation}\label{eq:shortcut-affine}
2^kT^k(n)=3^sn+C(w),
\end{equation}
where $C(w)=C_k$.
If the one-bits of $w$ occur at positions
\[
0\le d_1<\cdots<d_s<k,
\]
then
\begin{equation}\label{eq:carry-one-positions}
C(w)=\sum_{j=1}^{s}2^{d_j}3^{s-j}.
\end{equation}
\end{lemma}

\begin{proof}
Equation \eqref{eq:shortcut-affine} follows by induction from
\eqref{eq:shortcut-map}.  If $\eps_j=0$, neither the power of $3$ nor the
carry changes.  If $\eps_j=1$, multiplication by $3$ sends the previous
carry to $3C_j$, and the new $+1$ contributes $2^j$ after clearing the
common denominator $2^{j+1}$.  Expanding the recurrence gives
\eqref{eq:carry-one-positions}.
\end{proof}

\begin{lemma}[Parity cylinder]\label{lem:parity-cylinder}
For every word $w\in\set{0,1}^k$ there is exactly one residue class modulo
$2^k$ whose elements have parity prefix $w$.  Its canonical representative
is
\begin{equation}\label{eq:rho-shortcut}
\rho(w)=\left[-3^{-s}C(w)\right]_{2^k},
\qquad 0\le\rho(w)<2^k,
\end{equation}
where $[\cdot]_{2^k}$ denotes the least nonnegative representative.
\end{lemma}

\begin{proof}
Reducing \eqref{eq:shortcut-affine} modulo $2^k$ gives the necessary
congruence.  Conversely, the standard parity-vector bijection identifies
$n\bmod2^k$ with the first $k$ parities of its orbit; equivalently, the
recursion may be read backwards one bit at a time.
\end{proof}

The multiplicative coefficient of a prefix is
\begin{equation}\label{eq:q-shortcut}
q_j(w)=\frac{3^{s_j}}{2^j}.
\end{equation}

\subsection{Accelerated valuation words}

On positive odd integers define the accelerated map
\begin{equation}\label{eq:U-map}
U(x)=\frac{3x+1}{2^{\vtwo(3x+1)}}.
\end{equation}
A valuation word is
\[
a=(a_0,\ldots,a_{r-1}),
\qquad a_i\ge1.
\]
Put
\[
A_0=0,
\qquad
A_j=\sum_{i=0}^{j-1}a_i,
\]
and define
\begin{equation}\label{eq:valuation-carry}
D_0=0,
\qquad
D_{j+1}=3D_j+2^{A_j}.
\end{equation}

\begin{lemma}[Affine valuation-word identity]\label{lem:valuation-affine}
If $a$ is the first $r$ accelerated valuations of an odd integer $x$, and
$X=U^r(x)$, then
\begin{equation}\label{eq:valuation-affine}
2^{A_r}X=3^rx+D(a).
\end{equation}
If $a=uv$, then
\begin{equation}\label{eq:valuation-concat}
D(uv)=3^{|v|}D(u)+2^{A(u)}D(v).
\end{equation}
\end{lemma}

\begin{proof}
Both statements follow immediately by induction and comparison of the
constant terms under composition.
\end{proof}

\section{Primitive paradoxes and near-resonance}\label{sec:paradox}

\begin{definition}[Paradoxical block]
Let $x_i=T^i(x_0)$, and let a block have length $k$ and $s$ odd inputs.  It
is paradoxical if
\[
\frac{3^s}{2^k}<1
\qquad\text{but}\qquad
x_k\ge x_0.
\]
It is \emph{primitive} if
\begin{equation}\label{eq:primitive-barrier}
\frac{3^{s_j}}{2^j}\ge1
\qquad(0\le j<k),
\end{equation}
so that $k$ is the first coefficient crossing below one.
\end{definition}

Put
\[
\alpha=\log_2 3.
\]

\begin{theorem}[Primitive paradox resonance]\label{thm:primitive-resonance}
Let $x_0,\ldots,x_k$ be a primitive paradoxical block, and suppose
\[
m=x_0=\min_{0\le i\le k}x_i.
\]
Let $s$ be the number of odd inputs, and let
\[
0\le d_1<\cdots<d_s<k
\]
be their positions.  Then
\begin{equation}\label{eq:k-floor-alpha-s}
k=\floor{\alpha s}+1,
\end{equation}
\begin{equation}\label{eq:odd-position-bound}
d_j\le\floor{\alpha(j-1)}
\qquad(1\le j\le s),
\end{equation}
and
\begin{equation}\label{eq:resonance-bound}
0<\frac{k}{s}-\alpha
<\frac{1}{3m\log2}.
\end{equation}
Moreover,
\begin{equation}\label{eq:minimum-carry-bound}
m
\le
\frac{\displaystyle \frac13
 \sum_{j=0}^{s-1}2^{-\{\alpha j\}}}
{\displaystyle 2^{1-\{\alpha s\}}-1}
<
\frac{s}
{3\log2\bigl(1-\{\alpha s\}\bigr)}.
\end{equation}
If
\begin{equation}\label{eq:legendre-threshold}
m>\frac{2s^2}{3\log2},
\end{equation}
then the reduced fraction associated with $k/s$ is an upper continued-fraction
convergent of $\alpha$.
\end{theorem}

\begin{proof}
At an odd input $x_i$,
\[
\frac{x_{i+1}}{x_i}
=
\frac32\left(1+\frac{1}{3x_i}\right),
\]
whereas an even input contributes $1/2$.  Therefore
\begin{equation}\label{eq:product-correction}
\frac{x_k}{x_0}
=
\frac{3^s}{2^k}
\prod_{\substack{0\le i<k\\x_i\text{ odd}}}
\left(1+\frac1{3x_i}\right).
\end{equation}
Set
\[
\delta=k\log2-s\log3.
\]
The final coefficient is below one, so $\delta>0$.  Since $x_k\ge x_0$ and
$x_i\ge m$,
\[
\delta
\le
\sum_{x_i\text{ odd}}
\log\left(1+\frac1{3x_i}\right)
\le
s\log\left(1+\frac1{3m}\right)
<\frac{s}{3m}.
\]
Dividing by $s\log2$ proves \eqref{eq:resonance-bound}.

The last step must be even; otherwise the coefficient would be multiplied by
$3/2$ and could not cross from at least one to below one.  Hence
\[
2^{k-1}\le3^s<2^k,
\]
which proves \eqref{eq:k-floor-alpha-s}.

Before the $j$th odd input, the block has length $d_j$ and has encountered
$j-1$ odd inputs.  The primitive condition gives
\[
2^{d_j}\le3^{j-1},
\]
which is \eqref{eq:odd-position-bound}.

By Lemma~\ref{lem:shortcut-affine},
\[
2^kx_k=3^sm+C(w).
\]
Since $x_k\ge m$,
\begin{equation}\label{eq:defect-carry}
(2^k-3^s)m\le C(w).
\end{equation}
Using \eqref{eq:odd-position-bound} in
\eqref{eq:carry-one-positions},
\[
\begin{aligned}
C(w)
&\le
\sum_{j=1}^s
2^{\floor{\alpha(j-1)}}3^{s-j}\\
&=
3^{s-1}\sum_{j=0}^{s-1}2^{-\{\alpha j\}}.
\end{aligned}
\]
Also, by \eqref{eq:k-floor-alpha-s},
\[
2^k-3^s
=
3^s\left(2^{1-\{\alpha s\}}-1\right).
\]
Combining these identities with \eqref{eq:defect-carry} gives the first
bound in \eqref{eq:minimum-carry-bound}.  The second follows from
\[
2^u-1\ge u\log2
\qquad(0\le u\le1)
\]
and the trivial estimate of the numerator by $s/3$.

Finally, \eqref{eq:resonance-bound} and
\eqref{eq:legendre-threshold} imply
\[
0<\frac{k}{s}-\alpha<\frac1{2s^2}.
\]
Legendre's continued-fraction criterion applies to the reduced fraction.
\end{proof}

\begin{proposition}[Harmonic correction bound]\label{prop:harmonic-correction}
Suppose the block in Theorem~\ref{thm:primitive-resonance} is injective.
Then
\begin{equation}\label{eq:harmonic-drift}
0<k\log2-s\log3
\le
\frac1{3m}
+
\frac16\log\left(1+\frac{2s}{m}\right).
\end{equation}
\end{proposition}

\begin{proof}
The odd inputs are distinct odd integers at least $m$.  Hence
\[
\sum_{x_i\text{ odd}}\frac1{x_i}
\le
\sum_{j=0}^{s-1}\frac1{m+2j}
\le
\frac1m+\frac12\log\left(1+\frac{2s}{m}\right).
\]
Apply $\log(1+u)\le u$ to \eqref{eq:product-correction}.
\end{proof}

\begin{remark}
Theorem~\ref{thm:primitive-resonance} explains why ``the starting value is
large, so the $+1$ is negligible'' is not by itself a proof.  A large
minimum forces an increasingly accurate near-resonance between $2^k$ and
$3^s$, and continued fractions provide infinitely many such rational
approximations.  The remaining information is arithmetic: the unique
$2$-adic seed residue must land in the short Archimedean interval allowed by
\eqref{eq:defect-carry}.
\end{remark}

\section{Inverse fibers and an exact reformulation of Collatz}\label{sec:fibers}

For an odd positive integer $n$, define
\begin{equation}\label{eq:Mr-def}
M_r(n)=
\min\set{y>0\text{ odd}:U^r(y)=U^r(n)}.
\end{equation}
The set is nonempty because it contains $n$.

\begin{lemma}[Nested inverse minima]\label{lem:nested-minima}
For every odd positive integer $n$,
\begin{equation}\label{eq:Mr-monotone}
M_{r+1}(n)\le M_r(n)\le n.
\end{equation}
\end{lemma}

\begin{proof}
The upper bound is immediate.  If $y$ realizes $M_r(n)$, append the actual
next valuation of the orbit of $n$ to the inverse word carrying $y$ to
$U^r(n)$.  The same starting value then reaches $U^{r+1}(n)$ in $r+1$
accelerated steps.  Thus $M_{r+1}(n)\le M_r(n)$.
\end{proof}

\begin{theorem}[Stationary inverse-fiber extinction]\label{thm:SFE-equivalence}
The following are equivalent.
\begin{enumerate}[label=(\roman*),leftmargin=2.3em]
\item Every positive integer reaches $1$ under the Collatz map.
\item For every odd $n>1$, there exists $r\ge1$ such that
      \begin{equation}\label{eq:SFE}
      M_r(n)<n.
      \end{equation}
\item No odd $n>1$ is the least positive member of every finite inverse fiber
      above its own forward orbit.
\end{enumerate}
\end{theorem}

\begin{proof}
Assume Collatz.  For an odd $n>1$, choose $r$ with $U^r(n)=1$.  Since
$U^r(1)=1$, one has $M_r(n)=1<n$.

Conversely, assume \eqref{eq:SFE}.  We prove by strong induction on odd $n$
that every odd $n$ reaches $1$.  The case $n=1$ is trivial.  Let $n>1$ and
assume every smaller positive odd integer reaches $1$.  Choose $r$ and
$0<y<n$ with
\[
U^r(y)=U^r(n).
\]
By induction, $y$ reaches $1$; hence the common endpoint reaches $1$, and so
does $n$.  Every positive even integer becomes odd after finitely many
halvings.  This proves Collatz.
\end{proof}

\begin{definition}[Stationary Bellman code]
An odd integer $n>1$ is a stationary Bellman code if
\begin{equation}\label{eq:stationary-Bellman}
M_r(n)=n
\qquad(r\ge0).
\end{equation}
\end{definition}

A least counterexample to Collatz would be a stationary Bellman code.  The
entire problem is therefore reduced to excluding one positive integer that
remains the least positive cost in every finite inverse tree.

\section{Terminal locality and the exact local--global score}\label{sec:local-global}

\begin{theorem}[Terminal $3$-adic locality]\label{thm:terminal-locality}
Let a valuation word be written $a=uv$, where $|v|=\ell$, and let $X$ be the
endpoint of any realization of $a$.  Then
\begin{equation}\label{eq:terminal-locality}
3^\ell\mid X+1
\quad\Longleftrightarrow\quad
3^\ell\mid D(v)+2^{A(v)}.
\end{equation}
Thus the condition $\vthree(X+1)\ge\ell$ depends only on the final $\ell$
valuations.
\end{theorem}

\begin{proof}
By Lemma~\ref{lem:valuation-affine}, reducing the endpoint identity modulo
$3^\ell$ gives
\[
3^\ell\mid X+1
\quad\Longleftrightarrow\quad
3^\ell\mid D(a)+2^{A(a)}.
\]
Using \eqref{eq:valuation-concat},
\[
D(a)+2^{A(a)}
=
3^\ell D(u)+2^{A(u)}\bigl(D(v)+2^{A(v)}\bigr).
\]
The factor $2^{A(u)}$ is invertible modulo $3^\ell$.
\end{proof}

Now let an actual word have the form
\[
a=u\tau,
\qquad
|u|=d,
\quad A(u)=P,
\quad |\tau|=k,
\quad A(\tau)=T.
\]
Let $R$ be the starting value and
\[
Z=U^d(R).
\]
Consider a competing word $b$ of length $k$ and total exponent $B$, and put
\[
c=T-B.
\]
Compatibility with the same endpoint is equivalent to
\begin{equation}\label{eq:local-compatibility}
3^k\mid D(\tau)-2^cD(b).
\end{equation}
Define the local center
\begin{equation}\label{eq:local-center}
h(\tau,b)
=
\frac{D(\tau)-2^cD(b)}{3^k}
\in\Z.
\end{equation}

\begin{proposition}[Local--global score]\label{prop:local-global-score}
Let $Y$ be the competing ancestor along $b$.  Then
\begin{equation}\label{eq:local-global-identity}
2^cY=Z+h(\tau,b).
\end{equation}
Consequently,
\begin{equation}\label{eq:score-identity}
Y-R
=
\frac{\Sigma(u,\tau,c)+h(\tau,b)}{2^c},
\end{equation}
where
\begin{equation}\label{eq:prefix-surplus}
\Sigma(u,\tau,c)=Z-2^cR.
\end{equation}
The compatibility term $h$ is determined by the terminal suffix and the
competitor, whereas $\Sigma$ depends on the complete prefix and the actual
starting value.
\end{proposition}

\begin{proof}
The actual suffix equation is
\[
2^TX=3^kZ+D(\tau),
\]
and the competing inverse equation is
\[
3^kY=2^BX-D(b).
\]
Multiplying the second equation by $2^c$, substituting the first, and using
\eqref{eq:local-center} gives \eqref{eq:local-global-identity}.  Subtracting
$R$ gives \eqref{eq:score-identity}.
\end{proof}

If a terminal fiber with fixed $B$ is nonempty, let
\[
h_{\max}(\tau,B)=\max_b h(\tau,b).
\]
Then a smaller positive compatible ancestor exists whenever
\begin{equation}\label{eq:score-negative}
-R2^c<\Sigma(u,\tau,c)+h_{\max}(\tau,B)<0.
\end{equation}
The right inequality is the descent condition; the left inequality is
positivity.  This is the exact synchronization score.

\begin{remark}
Terminal locality does not settle the sign of
\eqref{eq:score-negative}.  The same terminal suffix and the same slope data
can be attached to different prefixes whose Archimedean surpluses have
opposite signs.  Any successful proof must retain both ends of the word.
\end{remark}

\section{The positive-integer image in parity coordinates}\label{sec:inverse-carry}

Let $\eps=(\eps_0,\eps_1,\ldots)$ be an infinite binary sequence, and put
\[
s_k=\sum_{j=0}^{k-1}\eps_j.
\]
For every prefix define
\begin{equation}\label{eq:rho-k}
\rho_k
=
\left[
-\sum_{j=0}^{k-1}
\frac{\eps_j2^j}{3^{s_{j+1}}}
\right]_{2^k},
\qquad
0\le\rho_k<2^k.
\end{equation}
The representatives are compatible, so there is a unique bit $b_k$ with
\begin{equation}\label{eq:inverse-carry-bit}
\rho_{k+1}=\rho_k+b_k2^k,
\qquad b_k\in\set{0,1}.
\end{equation}

Let $C_k$ be the shortcut carry from \eqref{eq:shortcut-carry-rec}, and
define the canonical endpoint
\begin{equation}\label{eq:z-k}
z_k=\frac{3^{s_k}\rho_k+C_k}{2^k}.
\end{equation}

\begin{proposition}[Intrinsic inverse-carry recursion]\label{prop:carry-recursion}
Starting from $\rho_0=z_0=s_0=0$, one has
\begin{equation}\label{eq:b-parity}
b_k\equiv z_k+\eps_k\pmod2,
\end{equation}
\begin{equation}\label{eq:z-recursion}
z_{k+1}
=
\frac{3^{\eps_k}z_k+\eps_k+b_k3^{s_{k+1}}}{2},
\end{equation}
and
\begin{equation}\label{eq:z-range}
0\le z_k<3^{s_k}.
\end{equation}
\end{proposition}

\begin{proof}
Replacing $\rho_k$ by $\rho_k+b2^k$ changes the endpoint after the first $k$
steps from $z_k$ to $z_k+b3^{s_k}$.  This value must have parity $\eps_k$.
Since $3^{s_k}$ is odd, this gives \eqref{eq:b-parity}.  Applying the next
branch gives \eqref{eq:z-recursion}.

The range follows by induction.  If $\eps_k=0$, the two possibilities are
$z_{k+1}=z_k/2$ or $(z_k+3^{s_k})/2$.  If $\eps_k=1$, they are
$(3z_k+1)/2$ or $(3z_k+1+3^{s_k+1})/2$; the parity restriction in each case
ensures the strict upper bound $3^{s_{k+1}}$.
\end{proof}

\begin{theorem}[Finite-support characterization]\label{thm:finite-support}
For an infinite parity sequence $\eps$, the following are equivalent.
\begin{enumerate}[label=(\roman*),leftmargin=2.3em]
\item The unique $2$-adic integer with parity sequence $\eps$ is a
      nonnegative ordinary integer.
\item $b_k=0$ for all sufficiently large $k$.
\item $\rho_k$ is eventually constant.
\item $\sup_k\rho_k<\infty$.
\item $\rho_k/2^k\to0$.
\end{enumerate}
\end{theorem}

\begin{proof}
Equation \eqref{eq:inverse-carry-bit} shows that $\rho_k$ is the partial
binary expansion
\[
\rho_k=\sum_{j=0}^{k-1}b_j2^j.
\]
Thus eventual vanishing of $b_j$, eventual constancy of $\rho_k$, and
boundedness of $\rho_k$ are equivalent.  Their common eventual value is the
corresponding nonnegative ordinary integer.

If $\rho_k$ is eventually constant then $\rho_k/2^k\to0$.  Conversely, if
$b_k=1$ infinitely often, then at each such stage
\[
\frac{\rho_{k+1}}{2^{k+1}}\ge\frac12,
\]
so the ratio cannot tend to zero.
\end{proof}

\subsection{The mixed dyadic--ternary corner}

Assume the coefficient barrier
\begin{equation}\label{eq:shortcut-barrier}
3^{s_k}\ge2^k
\qquad(k\ge0).
\end{equation}
Put
\begin{equation}\label{eq:Rk-carry}
R_k=\frac{C_k}{3^{s_k}}
=
\sum_{j=0}^{k-1}
\frac{\eps_j2^j}{3^{s_{j+1}}}.
\end{equation}
From \eqref{eq:z-k},
\begin{equation}\label{eq:corner-identity}
\frac{z_k}{3^{s_k}}
=
\frac{\rho_k+R_k}{2^k}.
\end{equation}
If $\eps_j=1$, the barrier gives
\[
\frac{2^j}{3^{s_{j+1}}}
=
\frac13\frac{2^j}{3^{s_j}}
\le\frac13.
\]
Therefore
\[
0\le R_k\le\frac{s_k}{3}\le\frac{k}{3}.
\]

\begin{corollary}[Corner criterion]\label{cor:corner-criterion}
Under the barrier \eqref{eq:shortcut-barrier},
\begin{equation}\label{eq:corner-bound}
0\le
\frac{z_k}{3^{s_k}}-\frac{\rho_k}{2^k}
\le
\frac{k}{3\cdot2^k}.
\end{equation}
Consequently, the parity sequence comes from a nonnegative ordinary integer
if and only if
\begin{equation}\label{eq:corner-limit}
\frac{z_k}{3^{s_k}}\longrightarrow0.
\end{equation}
\end{corollary}

\subsection{Divergent orbits and high surplus}

The following reduction uses the orbit-sparsity theorem of Garcia and Tal
\cite{GarciaTal1999}.

\begin{proposition}[Barrier normalization of divergence]\label{prop:barrier-normalization}
If a positive shortcut Collatz orbit is nonperiodic and does not reach the
trivial cycle, then it tends to infinity.  Along such an orbit,
\[
\sum_{j\ge0}\frac1{x_j}<\infty.
\]
Consequently, if the orbit is divergent, its coefficient sequence has a
global minimum; after shifting to a time of global minimum, the relative
coefficient sequence satisfies
\begin{equation}\label{eq:shifted-barrier}
q_k\ge1
\qquad(k\ge0),
\end{equation}
and
\begin{equation}\label{eq:q-to-infty}
q_k\longrightarrow\infty.
\end{equation}
\end{proposition}

\begin{proof}
A nonperiodic deterministic orbit has no repeated value, hence eventually
leaves every finite interval.  Garcia and Tal prove a power-saving upper
bound for the number of points of a divergent orbit in an interval.  Summing
that estimate over dyadic shells gives convergence of the reciprocal sum.

As in \eqref{eq:product-correction},
\[
\frac{x_k}{x_0}
=q_k
\prod_{\substack{0\le j<k\\x_j\text{ odd}}}
\left(1+\frac1{3x_j}\right).
\]
The reciprocal-sum convergence implies convergence of the product to a
finite positive limit.  Since $x_k\to\infty$, one obtains $q_k\to\infty$.
The positive sequence $q_k$ therefore attains a global minimum.  Shifting to
a minimizing time gives \eqref{eq:shifted-barrier} and
\eqref{eq:q-to-infty}.
\end{proof}

Put
\begin{equation}\label{eq:theta}
\theta=\log_3 2
\end{equation}
and define the critical surplus
\begin{equation}\label{eq:critical-surplus}
e_k=s_k-\ceil{\theta k}.
\end{equation}
Under the barrier, $e_k\ge0$, and
\begin{equation}\label{eq:q-surplus}
q_k
=
3^{e_k}
\frac{3^{\ceil{\theta k}}}{2^k},
\qquad
1\le
\frac{3^{\ceil{\theta k}}}{2^k}
<3.
\end{equation}

\begin{corollary}[High-surplus necessity]\label{cor:high-surplus}
A divergent orbit shifted into the barrier satisfies
\begin{equation}\label{eq:e-to-infty}
e_k\longrightarrow\infty.
\end{equation}
At the same time, because the parity sequence comes from one positive
ordinary integer, its inverse-carry bits satisfy
\begin{equation}\label{eq:b-eventual-zero}
b_k=0
\quad\text{for all sufficiently large }k.
\end{equation}
\end{corollary}

Thus the divergence half of Collatz is reduced to excluding the coexistence
\[
e_k\to\infty
\qquad\text{and}\qquad
b_k=0\text{ eventually}.
\]

\section{Transposition winding and mechanical extremizers}\label{sec:winding}

The unordered-majorization order used by Rozier and Terracol moves one-bits
left or right by adjacent transpositions.  Real affine carries respond
monotonically to these moves, but the canonical $2$-adic seed can wrap around
its modulus.  The following identity records that wrap exactly.

\begin{lemma}[Adjacent-transposition winding]\label{lem:winding}
Let
\[
w=u10v,
\qquad
w'=u01v,
\]
where
\[
|u|=i,
\qquad
\wt(u)=p,
\qquad
\wt(v)=t.
\]
Let both words have total length $k$ and weight
\[
s=p+t+1.
\]
Then
\begin{equation}\label{eq:carry-transposition}
C(w')-C(w)=2^i3^t.
\end{equation}
Let $a$ be the least positive odd representative satisfying
\begin{equation}\label{eq:a-transposition}
a\equiv3^{-(p+1)}\pmod{2^{k-i}},
\qquad
1\le a<2^{k-i}.
\end{equation}
Then
\begin{equation}\label{eq:rho-transposition}
\rho(w')
=
\rho(w)-2^ia+\eta2^k,
\qquad
\eta\in\set{0,1},
\end{equation}
where
\begin{equation}\label{eq:eta-wrap}
\eta=1
\quad\Longleftrightarrow\quad
\rho(w)<2^ia.
\end{equation}
\end{lemma}

\begin{proof}
The moved one contributes $2^i3^t$ in $w$ and $2^{i+1}3^t$ in $w'$, while
all other contributions are unchanged.  This proves
\eqref{eq:carry-transposition}.

By Lemma~\ref{lem:parity-cylinder},
\[
\rho(w')-\rho(w)
\equiv
-3^{-s}\bigl(C(w')-C(w)\bigr)
\equiv
-2^i3^{-(p+1)}
\pmod{2^k}.
\]
Factoring out $2^i$ gives \eqref{eq:rho-transposition}.  Since both
representatives lie in $[0,2^k)$, at most one wrap by $2^k$ occurs, and the
criterion is \eqref{eq:eta-wrap}.
\end{proof}

Define the endpoint of the canonical realization by
\[
Z(w)=T^k(\rho(w))
\]
and the first-crossing descent margin
\begin{equation}\label{eq:J-def}
J(w)=\bigl(2^k-3^s\bigr)\rho(w)-C(w).
\end{equation}
Then
\begin{equation}\label{eq:J-endpoint}
J(w)=2^k\bigl(\rho(w)-Z(w)\bigr).
\end{equation}
Hence $J(w)>0$ is exactly descent of the least positive realization.

\begin{corollary}[Winding cocycle for descent]\label{cor:winding-cocycle}
Under the transposition in Lemma~\ref{lem:winding},
\begin{equation}\label{eq:J-transposition}
J(w')-J(w)
=
\bigl(2^k-3^s\bigr)\bigl(-2^ia+\eta2^k\bigr)
-2^i3^t.
\end{equation}
\end{corollary}

The deterministic carry term is favorable, but the binary winding term
$\eta2^k$ can reverse its effect.  This is the finite version of the
local--global obstruction.

\subsection{Barrier-preserving reduction to the latest-one word}

Let a primitive first-crossing word have one-positions
\[
d_1<\cdots<d_s.
\]
Define
\begin{equation}\label{eq:latest-one-positions}
D_j=\floor{\alpha(j-1)},
\qquad
\alpha=\log_2 3.
\end{equation}

\begin{theorem}[Barrier-preserving transposition chain]\label{thm:mechanical-chain}
Every primitive first-crossing word can be transformed into the word whose
one-positions are $D_1,\ldots,D_s$ by finitely many adjacent moves
\[
10\longmapsto01.
\]
Every intermediate word remains a primitive first-crossing word of the same
length and weight.
\end{theorem}

\begin{proof}
The position bound \eqref{eq:odd-position-bound} gives $d_j\le D_j$.  Move
the ones from right to left.  Suppose the $j$th one is moved from position
$i$ to $i+1$, where $i+1\le D_j$.  The only prefix whose one-count is
reduced is the prefix of length $i+1$; after the move it contains $j-1$
ones.  Since
\[
i+1\le D_j\le\alpha(j-1),
\]
one has
\[
2^{i+1}\le3^{j-1}.
\]
Thus the coefficient barrier remains valid.  The total length and total
weight are unchanged, so the final crossing remains unchanged.
\end{proof}

Theorem~\ref{thm:mechanical-chain} and Corollary~\ref{cor:winding-cocycle}
reduce first-crossing descent to an accumulated winding estimate along a
finite path to a mechanical extremizer.  The real carry order is understood;
the remaining datum is the sequence of wrap bits $\eta$.

\section{Equal-slope carry returns}\label{sec:equal-slope}

For a shortcut word $w$ of length $k$ and weight $s$, define its equal-slope
carry spectrum
\begin{equation}\label{eq:carry-spectrum}
\mathcal K_{k,s}(w)
=
\set{
\frac{C(v)-C(w)}{3^s}:
\begin{array}{l}
|v|=k,\ \wt(v)=s,\\
C(v)\equiv C(w)\pmod{3^s}
\end{array}
}.
\end{equation}

\begin{theorem}[Equal-slope coalescence]\label{thm:equal-slope}
Let $w$ be the actual first $k$ parity bits of a positive integer $n$, let
$s=\wt(w)$, and put $X=T^k(n)$.  Suppose there is a word $v$ of length $k$
and weight $s$, and an integer $d$ satisfying
\begin{equation}\label{eq:equal-slope-difference}
C(v)=C(w)+d3^s,
\qquad
1\le d<n.
\end{equation}
Then $y=n-d$ is positive, its first $k$ parity bits are $v$, and
\begin{equation}\label{eq:equal-slope-coalesce}
T^k(y)=X.
\end{equation}
\end{theorem}

\begin{proof}
Using \eqref{eq:equal-slope-difference},
\[
3^sy+C(v)
=3^s(n-d)+C(w)+d3^s
=3^sn+C(w)
=2^kX.
\]
Thus $y$ satisfies the canonical congruence for the word $v$ and reaches the
same endpoint.  Lemma~\ref{lem:parity-cylinder} gives exactness of the parity
prefix.
\end{proof}

\begin{corollary}[Equal-slope stationary extinction criterion]
If every nonterminating positive orbit beginning at $n>1$ has an actual
prefix $w$ for which
\begin{equation}\label{eq:equal-slope-spectrum-target}
\mathcal K_{k,s}(w)\cap\set{1,\ldots,n-1}\ne\varnothing,
\end{equation}
then Collatz holds.
\end{corollary}

\begin{proof}
Theorem~\ref{thm:equal-slope} supplies a smaller positive coalescing ancestor;
Theorem~\ref{thm:SFE-equivalence} finishes the proof.
\end{proof}

Unlike an asymptotic slope comparison, the criterion
\eqref{eq:equal-slope-difference} is pointwise and integer-exact.  Its
remaining difficulty is recurrence: such a carry return need not occur at
every depth, and the desired theorem must force one occurrence along every
positive stationary ray.

\section{Ten research programs for synchronization}\label{sec:programs}

This section records ten routes to stationary inverse-fiber extinction.  Each
subsection states an exact target strong enough to imply a full proof, then
identifies the obstruction left by current methods.

\subsection{Program I: winding coercivity along mechanical chains}

For a barrier-preserving chain
\[
w=w_0\to w_1\to\cdots\to w_N=w_{\mathrm{mech}},
\]
let $\eta_j$ be the winding bit of the $j$th transposition.  Iterating
\eqref{eq:J-transposition} gives an exact additive cocycle for $J(w_N)-J(w)$.

\begin{problem}[Winding-coercivity theorem]\label{prob:winding-coercivity}
Prove that every nontrivial primitive first-crossing word $w$ realized by the
stabilized residue of one positive ordinary integer satisfies
\begin{equation}\label{eq:winding-coercive-target}
J(w)>0.
\end{equation}
Equivalently, prove that the cumulative positive wrap contribution cannot
overcome the deterministic nonwrap and carry losses in
\eqref{eq:J-transposition}.
\end{problem}

A proof would establish first-crossing descent and hence Terras's
coefficient-stopping-time assertion.  The central subproblem is to show that
large-weight transpositions cannot wrap too often on a stationary
positive-integer ray.  Mechanical and Christoffel extremizers supply the
correct endpoint of the chain \cite{FernandezIbanez2026}; the new arithmetic
content must control the winding bits.

\subsection{Program II: equal-slope carry recurrence}

For the actual prefix $w_k(n)$ define
\[
\mathcal K_k(n)=\mathcal K_{k,s_k}(w_k(n)).
\]

\begin{problem}[Equal-slope carry recurrence]\label{prob:equal-slope-recurrence}
Prove that every nonterminating odd orbit beginning at $n>1$ has some depth
$k$ for which
\begin{equation}\label{eq:equal-slope-recurrence}
\mathcal K_k(n)\cap\set{1,\ldots,n-1}\ne\varnothing.
\end{equation}
\end{problem}

By Theorem~\ref{thm:equal-slope}, this is a direct proof of stationary
inverse-fiber extinction.  Adjacent transpositions generate carry changes of
the form $2^i3^t$; one seeks a structured signed combination equal to
$d3^s$ with $1\le d<n$.  The divisibility is triangular and $3$-adic, so a
promising organization is to choose swaps from the terminal end and correct
one ternary digit at a time.

Ordinary sum--product expansion does not directly apply: a single survivor
ray is graph-like and each fixed branch is $3$-adically contracting.  The
missing result is ray-dependent recurrence, not uniform expansion of the
whole residue language.

\subsection{Program III: Ostrowski surplus return sections}

Let
\[
e_k=s_k-\ceil{\theta k},
\qquad \theta=\log_3 2,
\]
and define the first activation time of surplus level $j$ by
\begin{equation}\label{eq:tau-j}
\tau_j=\min\set{k:e_k\ge j}.
\end{equation}
The blocks $[\tau_j,\tau_{j+1})$ are surplus return blocks over the Sturmian
critical scaffold of slope $\theta$.

\begin{problem}[Surplus-block carry innovation]\label{prob:surplus-innovation}
Prove that there is an absolute $L$ such that, for every barrier sequence and
every sufficiently large $j$, either
\begin{equation}\label{eq:surplus-carry-window}
b_k=1
\quad\text{for some }\tau_j\le k<\tau_{j+1}+L,
\end{equation}
or the corresponding marked surplus block belongs to a finite exceptional
family whose infinite concatenations represent eventually periodic negative
$2$-adic seeds.
\end{problem}

A divergent barrier orbit has $e_k\to\infty$ by
Corollary~\ref{cor:high-surplus}.  Thus repeated innovation would contradict
eventual vanishing of the $b_k$.  The continued fraction of $\theta$ and
Ostrowski numeration provide canonical return blocks; the desired theorem
must show that each new unit of surplus consumes a new high binary carry,
unless the code collapses into a periodic nonpositive sector.

\subsection{Program IV: a corner-escape Lyapunov function}

Normalize
\begin{equation}\label{eq:normalized-corner-state}
u_k=\frac{\rho_k}{2^k},
\qquad
v_k=\frac{z_k}{3^{s_k}},
\qquad
r_k=\frac{2^k}{3^{s_k}}.
\end{equation}
Then
\begin{equation}\label{eq:u-dynamics}
u_{k+1}=\frac{u_k+b_k}{2},
\end{equation}
\begin{equation}\label{eq:v-dynamics}
v_{k+1}
=
\frac{v_k+b_k}{2}
+
\frac{\eps_k}{2\cdot3^{s_{k+1}}},
\end{equation}
and
\begin{equation}\label{eq:r-dynamics}
r_{k+1}=\frac{2r_k}{3^{\eps_k}}.
\end{equation}
A divergent positive barrier orbit would approach the forbidden corner
\[
(u_k,v_k,r_k)\to(0,0,0)
\]
with $b_k=0$ eventually.

\begin{problem}[Corner-escape theorem]\label{prob:corner-escape}
Construct a coercive function
\[
\mathcal L(u,v,r,e,\phi),
\qquad
\phi=\{\theta k\},
\]
and constants $L,\delta>0$ such that every admissible length-$L$ block with
$b=0$ and $r\le1$ either violates the barrier or satisfies
\begin{equation}\label{eq:Lyapunov-increase}
\mathcal L_{k+L}\ge\mathcal L_k+\delta.
\end{equation}
The function must be bounded above on trajectories approaching
$(0,0,0)$.
\end{problem}

No one-step linear Lyapunov function is expected: the irrational phase of
the critical line is essential.  A successful function must retain an
unbounded state such as surplus, Ostrowski phase, or min-plus exposure gap;
bounded-memory potentials cannot distinguish a stationary ray from
terminating finite shadows.

\subsection{Program V: two-place separation of the plus orbit and minus shadow}

For a divergent positive barrier orbit, the same parity labels generate a
positive real ``minus'' shadow
\begin{equation}\label{eq:minus-shadow}
a_{t+1}
=
\begin{cases}
a_t/2,&\eps_t=0,\\[1mm]
(3a_t-1)/2,&\eps_t=1,
\end{cases}
\end{equation}
while the integer orbit satisfies
\begin{equation}\label{eq:plus-orbit}
x_{t+1}
=
\begin{cases}
x_t/2,&\eps_t=0,\\[1mm]
(3x_t+1)/2,&\eps_t=1.
\end{cases}
\end{equation}
The shadow series is positive in $\R$ and represents $-x_t$ in $\Q_2$.
Moreover $a_t/x_t\to0$.

\begin{lemma}[Pure multiplicative sum]\label{lem:pure-sum}
If $S_t=x_t+a_t$, then
\begin{equation}\label{eq:pure-sum-dynamics}
S_{t+1}
=
\begin{cases}
S_t/2,&\eps_t=0,\\[1mm]
3S_t/2,&\eps_t=1,
\end{cases}
\end{equation}
so $S_t=q_tS_0$.
\end{lemma}

\begin{proof}
The $+1$ and $-1$ cancel on odd steps; the even case is immediate.
\end{proof}

\begin{problem}[Two-place shadow separation]\label{prob:two-place-separation}
Prove that no positive integer trajectory $x_t$, positive real trajectory
$a_t>1$, and common label sequence $\eps_t$ can satisfy
\eqref{eq:minus-shadow}--\eqref{eq:plus-orbit},
\[
q_t\ge1,
\qquad q_t\to\infty,
\qquad \frac{a_t}{x_t}\to0,
\]
while every finite truncation of the shadow remains congruent to $-x_t$
modulo the corresponding power of two.
\end{problem}

This is an adelic separation problem.  Standard product-formula arguments do
not suffice because rational sequences may converge to unrelated limits at
different places.  One must derive additional algebraicity, recurrence, or
height control from the shared parity labels.

\subsection{Program VI: tropical curvature of the inverse-shadow envelope}

Fix an endpoint parameterization $X=3^rs-y$.  Every positive accelerated
$3x-1$ shadow word $v$ gives an affine line
\begin{equation}\label{eq:shadow-line}
\ell_v(s)=2^{A(v)}s-\mu_0(v),
\end{equation}
whose positive values are exactly the positive inverse ancestors of $X$.
Define the positive lower envelope
\begin{equation}\label{eq:lower-envelope}
E_r(s)=\min\set{\ell_v(s):\ell_v(s)>0,\ |v|=r}.
\end{equation}
At the actual parameter $s_r$,
\[
E_r(s_r)=M_r(n).
\]

\begin{problem}[Tropical exposure extinction]\label{prob:tropical-extinction}
Prove that no fixed positive line value $n>1$ can remain the exposed positive
minimum at the actual nested parameters for every depth.  Equivalently, if
$w_r$ is the actual inverse word, prove that for some $r$ there is a
competitor $v$ with
\begin{equation}\label{eq:tropical-target}
0<\ell_v(s_r)<\ell_{w_r}(s_r)=n.
\end{equation}
\end{problem}

This is exactly stationary inverse-fiber extinction in min-plus coordinates.
The hoped-for mechanism is cumulative curvature: an expanding family of
lines with different dyadic slopes should force a crossing near the actual
parameter.  Finite examples with large descent cutoffs show that no uniform
one-scale curvature constant exists; the theorem must be seed-dependent and
infinitary.

\subsection{Program VII: a moving-target Subspace Theorem}

Nested accelerated blocks give identities
\begin{equation}\label{eq:nested-trajectory-identity}
2^{A_{j,\ell}}x_{t_j+\ell}
-3^\ell x_{t_j}
=C_{j,\ell}.
\end{equation}
The carry is itself a sum of mixed powers of two and three.  A stationary
minimum supplies infinitely many related identities along one orbit.
Moving-target versions of Schmidt's Subspace Theorem allow the defining
forms to vary provided their height is small compared with the points
\cite{ZhaoLi2026}.

\begin{problem}[Nested small-form amplification]\label{prob:subspace-amplification}
Show that stationary inverse-fiber minimality forces infinitely many block
vectors $\mathbf X_j$ and moving linear forms $L_{v,i}^{(j)}$ satisfying
\begin{equation}\label{eq:subspace-inequality}
\prod_{v,i}
\frac{\abs{L_{v,i}^{(j)}(\mathbf X_j)}_v}
{\|\mathbf X_j\|_v}
\le H(\mathbf X_j)^{-\delta}
\end{equation}
for a fixed $\delta>0$, with target-form height
$o(h(\mathbf X_j))$.  Classify the resulting rational subspaces and show
that each forces either periodicity or a smaller coalescing ancestor.
\end{problem}

The missing preliminary step is amplification.  One trajectory identity is
exact but not unusually small relative to all of its coordinates; the proof
must manufacture several independent small forms from repeated near-returns,
marked carries, or nested endpoint comparisons.

\subsection{Program VIII: return-block $p$-adic interpolation}

A fixed block $w$ acts affinely on its realization class:
\begin{equation}\label{eq:block-affine-map}
F_w(x)=\lambda_wx+\beta_w.
\end{equation}
If the same marked block recurs, then
\begin{equation}\label{eq:block-iterate}
F_w^m(x)
=
\lambda_w^m x
+
\beta_w\frac{\lambda_w^m-1}{\lambda_w-1}.
\end{equation}
For primes $p\nmid6$, suitable powers of $\lambda_w$ may be close to one
$p$-adically, allowing interpolation of $m\mapsto F_w^m(x)$.

\begin{problem}[Marked recurrent-block theorem]\label{prob:padic-recurrence}
Prove that every positive stationary barrier ray contains one marked block
\[
(w,\rho\bmod2^a,z\bmod3^b,e\bmod m)
\]
recurring infinitely often in a form governed by a single $p$-adically
interpolable affine map.  Show that stationary minimality forces an
associated analytic descent function to have infinitely many zeros or a
fixed sign at infinitely many integers; use $p$-adic zero-counting to force
an identity, and classify the identity as periodic or descending.
\end{problem}

The obstacle is recurrence in the full marked state, not interpolation
itself.  Symbolic recurrence alone does not synchronize the seed residue,
endpoint phase, surplus level, and Bellman exposure status.

\subsection{Program IX: a gapped cross-prime fractal uncertainty principle}

A direct dyadic seed cylinder and ternary endpoint cylinder are exact affine
images, so naive independence is false.  Insert a middle gap of length $g$
and sum over admissible middle words.  This gives a transfer operator
\[
K_g:\ell^2(S_a)\longrightarrow\ell^2(E_b)
\]
between conditioned seed and endpoint residue sets.

\begin{problem}[Gapped arithmetic FUP]\label{prob:FUP}
After removing periodic structured sectors, prove constants $c>0$ and
$g_0$ for which
\begin{equation}\label{eq:FUP-target}
\left\|
\one_{E_b}K_g\one_{S_a}
\right\|_{2\to2}
\le e^{-cg}
\qquad(g\ge g_0)
\end{equation}
uniformly in the boundary data arising from positive stationary rays.
\end{problem}

An iterated stationary atom would retain unit mass through every scale,
contradicting \eqref{eq:FUP-target}.  Discrete Cantor-set uncertainty theory
suggests multiscale porosity and the exclusion of distributed spectral pairs
as the relevant hypotheses \cite{LaiShi2025}.  The unresolved Collatz input
is component-sensitive amplification: one exceptional ray must generate a
large enough family of neighboring competitors for an uncertainty estimate
to act on it.

\subsection{Program X: autoconjugacy, complementary drift, and Cobham rigidity}

The nontrivial $2$-adic autoconjugacy of Monks and Yazinski complements parity
bits \cite{MonksYazinski2004}.  A high-density barrier path is paired with a
low-density path having strongly contracting real drift, but its paired
$2$-adic seed need not be a positive ordinary integer.

\begin{problem}[Complement-and-repair synchronization]\label{prob:autoconjugacy-repair}
For every positive stationary barrier seed $n>1$, prove that some complemented
prefix can be modified by a terminal $3$-adically compatible correction to
produce a positive integer $y<n$ whose orbit reaches the actual endpoint of
$n$.
\end{problem}

One possible route is to prove that the original and complementary carry
streams admit finite-state or substitutive descriptions in multiplicatively
independent numeration systems.  Cobham-type rigidity would then force
eventual periodicity \cite{BustosFokkinkYassawi2025}.  Substitution cocycle
methods may provide a parallel Lyapunov criterion \cite{Solomyak2025}.  The
present gap is prior to those theorems: neither automaticity nor strong
irreducibility of the relevant joint carry cocycle has been established.

\section{Conditional completion theorems}\label{sec:completion}

Theorem~\ref{thm:SFE-equivalence} already gives the narrowest one-statement
completion criterion.  We record a second architecture that separates
finite coefficient crossings from infinite barrier rays.

\begin{conjecture}[First-crossing descent]\label{conj:FCD}
For every positive integer $n>1$, if $k$ is the first time for which
\[
q_k(n)<1,
\]
then
\begin{equation}\label{eq:FCD}
T^k(n)<n.
\end{equation}
\end{conjecture}

This is Terras's coefficient-stopping-time assertion in the present notation.

\begin{conjecture}[High-surplus carry innovation]\label{conj:HSCI}
Every infinite binary sequence satisfying
\begin{equation}\label{eq:HSCI-hyp}
s_k\ge\ceil{k\log_3 2}
\qquad(k\ge0)
\end{equation}
and
\begin{equation}\label{eq:HSCI-surplus}
s_k-\ceil{k\log_3 2}\longrightarrow\infty
\end{equation}
has infinitely many inverse-carry bits $b_k=1$.
\end{conjecture}

\begin{theorem}[Two-lemma conditional proof of Collatz]\label{thm:FCD-HSCI}
Conjectures~\ref{conj:FCD} and \ref{conj:HSCI} imply the Collatz conjecture.
\end{theorem}

\begin{proof}
Suppose Collatz is false and let $n>1$ be the least positive integer whose
orbit does not reach the trivial cycle.

If $q_k(n)<1$ for some $k$, choose the first such $k$.  By
Conjecture~\ref{conj:FCD},
\[
T^k(n)<n.
\]
The smaller integer lies on the nonterminating orbit of $n$, contradicting
minimality.

It remains that $q_k(n)\ge1$ for every $k$.  A bounded nonterminating orbit is
periodic.  But for a positive cycle, the affine cycle equation has positive
carry, hence the coefficient over one full period is strictly below one;
this would put us in the first case.  Thus the orbit of $n$ is divergent.

By Proposition~\ref{prop:barrier-normalization} and
Corollary~\ref{cor:high-surplus}, its critical surplus tends to infinity.
Since its parity sequence comes from the fixed positive integer $n$,
Theorem~\ref{thm:finite-support} says that $b_k=0$ eventually.  This
contradicts Conjecture~\ref{conj:HSCI}.
\end{proof}

\begin{remark}
The two conjectures divide the full problem cleanly.  First-crossing descent
handles every orbit whose coefficient ever falls below one, including every
positive cycle.  High-surplus carry innovation handles the remaining
infinite coefficient-barrier trajectories.  Stationary inverse-fiber
extinction is logically narrower as a single statement because it treats
both cases by one finite coalescence certificate.
\end{remark}

\section{What a successful synchronization proof must retain}

The preceding analysis gives several no-go principles.

\begin{enumerate}[label=(\roman*),leftmargin=2.5em]
\item \textbf{Finite-prefix recognition is insufficient.}
Every finite parity word is realized by positive integers, and can be
extended to a terminating orbit.  A counterexample can only be excluded by
using nestedness for one fixed seed.

\item \textbf{Terminal $3$-adic compatibility is insufficient.}
The suffix decides whether a competitor exists, but the prefix surplus
decides whether it is smaller and positive.

\item \textbf{Real drift is insufficient.}
Near-resonances between powers of two and three permit long paradoxical
behavior, as quantified by Theorem~\ref{thm:primitive-resonance}.

\item \textbf{Average entropy is insufficient.}
An exponentially sparse language may still contain one infinite stationary
ray.

\item \textbf{Naive cross-prime independence is false.}
Seed and endpoint cylinders are linked by exact affine identities.

\item \textbf{Ordinary approximate spectral gaps are insufficient.}
Arbitrarily long terminating shadows reproduce every finite parity pattern;
a norm supremum can choose a different seed at each depth and forget
stationarity.
\end{enumerate}

A successful invariant must therefore retain three pieces simultaneously:

\begin{equation}\label{eq:three-ends}
\boxed{
\begin{array}{c}
\text{left-end }2\text{-adic seed stabilization},\\
\text{right-end }3\text{-adic endpoint compatibility},\\
\text{Archimedean order between competing starts}.
\end{array}}
\end{equation}

The local--global score, transposition winding cocycle, equal-slope carry
spectrum, mixed corner dynamics, and tropical lower envelope are five
coordinate systems for this same three-part obstruction.

\section{Conclusion}

The present paper does not prove the Collatz conjecture.  It replaces a broad
request for ``global synchronization'' by one exact finite certificate:

\begin{equation}\label{eq:final-target}
\boxed{
\forall n>1\text{ odd}\quad
\exists r\ge1\quad
M_r(n)<n.
}
\end{equation}

By Theorem~\ref{thm:SFE-equivalence}, \eqref{eq:final-target} is equivalent
to Collatz.  It says that no positive integer other than $1$ can remain the
least positive point in every finite inverse fiber above its own orbit.

On the divergent branch, the obstruction becomes even more concrete.  A
counterexample shifted into the coefficient barrier would have to satisfy
simultaneously
\[
e_k\to\infty
\qquad\text{and}\qquad
b_k=0\text{ eventually}.
\]
Thus unbounded surplus above the critical mechanical word would coexist with
a finite-support inverse-carry stream.  Conjecture~\ref{conj:HSCI} asserts
that this is impossible.

The most direct arithmetic certificate found here is the equal-slope carry
return
\[
C(v)-C(w)=d3^s,
\qquad
1\le d<n,
\]
which immediately produces the smaller coalescing ancestor $n-d$.  The most
direct combinatorial obstruction is the winding bit in
\[
\rho(u01v)
=
\rho(u10v)-2^ia+\eta2^k.
\]
Christoffel and mechanical extremality control the real carry; controlling
the accumulated wraps would control descent.

The remaining theorem is therefore not a generic entropy, compactness, or
finite-state statement.  It is a cross-prime, two-ended, finite-input
rigidity theorem for one stationary arithmetic ray.

\begin{thebibliography}{99}

\bibitem{Barina2025}
D.~Barina,
\emph{Improved verification limit for the convergence of the Collatz conjecture},
J. Supercomput. \textbf{81} (2025), Article 810.

\bibitem{BustosFokkinkYassawi2025}
\'{A}.~Bustos-Gajardo, R.~Fokkink, and R.~Yassawi,
\emph{Cobham's theorem for the Gaussian integers},
arXiv:2510.01440, 2025.

\bibitem{FernandezIbanez2026}
C.~Fern\'andez and S.~Ib\'a\~nez,
\emph{Christoffel words as extremal structures in Collatz dynamics},
arXiv:2607.24844, 2026.

\bibitem{GarciaTal1999}
M.~V.~P. Garcia and F.~A. Tal,
\emph{A note on the generalized $3n+1$ problem},
Acta Arith. \textbf{90} (1999), no.~3, 245--250.

\bibitem{Kramer2026}
O.~Kramer,
\emph{Adaptive search in Collatz exponent-code space via $2$-adic and $3$-adic constraints},
arXiv:2607.10041, 2026.

\bibitem{Lagarias1985}
J.~C. Lagarias,
\emph{The $3x+1$ problem and its generalizations},
Amer. Math. Monthly \textbf{92} (1985), no.~1, 3--23.

\bibitem{LaiShi2025}
C.-K.~Lai and R.~Shi,
\emph{When is the fractal uncertainty principle for discrete Cantor sets most uncertain?},
arXiv:2501.00864, 2025.

\bibitem{Monks2025}
K.~G. Monks,
\emph{On $q$-analogs of the $3x+1$ dynamical system},
J. Integer Seq. \textbf{28} (2025), Article 25.6.8.

\bibitem{MonksYazinski2004}
K.~G. Monks and J.~Yazinski,
\emph{The autoconjugacy of the $3x+1$ function},
Discrete Math. \textbf{275} (2004), 219--236.

\bibitem{Niu2026}
T.~Niu,
\emph{Parity vectors and paradoxical sequences in the accelerated Collatz map},
arXiv:2605.13886, 2026.

\bibitem{RozierTerracol2025}
O.~Rozier and C.~Terracol,
\emph{Paradoxical behavior in Collatz sequences},
arXiv:2502.00948v5, 2025.

\bibitem{Solomyak2025}
B.~Solomyak,
\emph{On the Lyapunov spectrum of the twisted cocycle for substitutions},
arXiv:2502.14308, 2025.

\bibitem{Tao2022}
T.~Tao,
\emph{Almost all orbits of the Collatz map attain almost bounded values},
Forum Math. Pi \textbf{10} (2022), e12.

\bibitem{Terras1976}
R.~Terras,
\emph{A stopping time problem on the positive integers},
Acta Arith. \textbf{30} (1976), no.~3, 241--252.

\bibitem{YolcuAaronsonHeule2023}
E.~Yolcu, S.~Aaronson, and M.~J.~H. Heule,
\emph{An automated approach to the Collatz conjecture},
J. Autom. Reason. \textbf{67} (2023), Article 15.

\bibitem{ZhaoLi2026}
G.~Zhao and Y.~Li,
\emph{On a theorem of Schmidt's subspace theory with moving targets},
arXiv:2603.01627, 2026.

\end{thebibliography}

\end{document}
